
This section introduces the IQC objects used to study stability and performance
of interconnected systems in a Hilbert space. The construction builds on the
classical IQC framework~\cite{megretski1997system} and its
extension to hard IQCs for infinite-dimensional
channels~\cite{10156335}. We establish the generalized
system interconnection, formalize stability and performance, and define the
uncertainty and performance IQCs, as described in \cite{lenssen2025muanalysissynthesisframeworkpartial}. These definitions provide the foundation for
the theory developed in sections~\ref{sec:primal-controller-analysis},~\ref{sec:duality}, and~\ref{sec:dual-controller-synthesis}.

\subsection{Description of uncertain systems in a Hilbert-space}
We consider operators that map external inputs to external outputs through a
feedback interconnection between a linear (infinite dimensional) operators $P$
and $\Delta$.
This interconnection structure is formally defined as follows
\begin{definition}[Generalized Plant]\label{def:uncertain-interconnection}
Given operators
$P_\Delta :
L_{\mrm{e},[a,b]}^{\mbf{n}_{\mrm{w}_\Delta}}
\to
L_{\mrm{e},[a,b]}^{\mbf{n}_{\mrm{z}_\Delta}}$,
$P_{\Delta\mrm{p}} :
L_{\mrm{e},[a,b]}^{\mbf{n}_{\mrm{w}_{\mrm{p}}}}
\to
L_{\mrm{e},[a,b]}^{\mbf{n}_{\mrm{z}_\Delta}}$,
$P_{\mrm{p}\Delta} :
L_{\mrm{e},[a,b]}^{\mbf{n}_{\mrm{w}_\Delta}}
\to
L_{\mrm{e},[a,b]}^{\mbf{n}_{\mrm{z}_{\mrm{p}}}}$,
$P_{\mrm{p}} :
L_{\mrm{e},[a,b]}^{\mbf{n}_{\mrm{w}_{\mrm{p}}}}
\to
L_{\mrm{e},[a,b]}^{\mbf{n}_{\mrm{z}_{\mrm{p}}}}$,
and $\Delta :
L_{\mrm{e},[a,b]}^{\mbf{n}_{\mrm{z}_\Delta}}
\to
L_{\mrm{e},[a,b]}^{\mbf{n}_{\mrm{w}_\Delta}}$, define
\begin{equation}\label{eq:uncertain-map}
\bmat{z_\Delta\\ z_{\mrm{p}}}
= P\bmat{w_\Delta\\ w_{\mrm{p}}} =
\bmat{
P_\Delta & P_{\Delta\mrm{p}} \\
P_{\mrm{p}\Delta} & P_{\mrm{p}}
}\bmat{w_\Delta\\ w_{\mrm{p}}},
\quad
w_\Delta=\Delta z_\Delta.
\end{equation}
for input $w_{\mrm{p}} \in L_{\mrm{e},[a,b]}^{\mbf{n}_{\mrm{w}_{\mrm{p}}}}$ and output
$z_{\mrm{p}}\in L_{\mrm{e},[a,b]}^{\mbf{n}_{\mrm{z}_{\mrm{p}}}}$, there exist internal signals
$z_\Delta \in L_{\mrm{e},[a,b]}^{\mbf{n}_{\mrm{z}_\Delta}}$ and
$w_\Delta \in L_{\mrm{e},[a,b]}^{\mbf{n}_{\mrm{w}_\Delta}}$
satisfying the {interconnected system} $z_{\mrm{p}}= (P\star\Delta)w_{\mrm{p}}$.
\end{definition}
The interconnection structure described above captures the feedback interaction
between the nominal system $P$ and the uncertainty $\Delta$, where $w_{\mrm{p}}$ and $z_{\mrm{p}}$
represent external inputs and performance outputs, respectively, while
$w_\Delta$ and $z_\Delta$ form the internal feedback loop through the
uncertainty.
Before analyzing stability and performance, we must first ensure the
well-posedness of the interconnection, such that solutions exist, are unique,
and depend causally on inputs. This is the requirement recorded
in~(\ref{ass:well-posedness}).
Having established the class of uncertain systems and their well-posedness, we
now develop a framework to guarantee robust stability and performance for the
class of systems defined in Definition~\ref{def:uncertain-interconnection}.



\subsection{Robust stability and performance}
Veenman, Scherer, and K\"oro\u{g}lu~\cite{veenman2016robust} describe both
uncertainty and performance specifications through quadratic constraints on
interconnection signals.
For PIE systems, this formulation also supports structured robustness and
performance analysis~\cite{10156335,lenssen2025muanalysissynthesisframeworkpartial}.
We now define what it means for the system to remain stable despite
uncertainties.
Robust stability requires that the system is well-posed for all admissible
uncertainties and that the input-output map remains bounded. 
This involves removal of the performance channels in
\eqref{eq:uncertain-map}, and
introducing auxiliary signals $z_0$ and $w_0$.
\begin{definition}
The interconnection
\begin{align}\label{eqn:rob-stab-interconnect}
z_\Delta = P_\Delta w_\Delta + z_0, \quad
w_\Delta = \Delta z_\Delta + w_0,
\end{align}
with $z_0, w_0 \in L_{\mrm{e},[a,b]}$, is said to be
\emph{robustly stable} if
\begin{enumerate}
\item The uncertain system is well-posed for all $\Delta\in\mbf{\Delta}$.
\item The map from $\bmat{z_0\\w_0}$ to
$\bmat{z_\Delta\\w_\Delta}$ is bounded on $L_{\mrm{e},[a,b]}$
for all $\Delta\in\mbf{\Delta}$.
\end{enumerate}
\end{definition}
Within the IQC framework, the uncertainty operator $\Delta$ is described by a
quadratic constraint defined by an appropriate multiplier. A stability
certificate combines this constraint with a complementary condition on the
plant. We use the finite-horizon, or hard, IQC
formulation~\cite{6915700,10156335}, leading to the following definition.
\begin{definition}\label{def:uncertainty-iqc}

An uncertainty operator
$\Delta:
L_{\mrm{e},[a,b]}^{\mbf{n}_{\mrm{z}_\Delta}}
\to
L_{\mrm{e},[a,b]}^{\mbf{n}_{\mrm{w}_\Delta}}$
is said to satisfy the hard IQC $(\Psi_\Delta, {\mcl V_\Delta})$ defined by the dynamic operator
$\Psi_\Delta:
L^{\mbf{n}_{\mrm{z}_\Delta} + \mbf{n}_{\mrm{w}_\Delta}}_{\mrm{e}, [a,b]}
\to
L^{\mbf{n}_{\tilde{\mrm{z}}_{\Delta}} + \mbf{n}_{\tilde{\mrm{w}}_{\Delta}}}_{\mrm{e}, [a,b]}$
and 
linear bounded $\mcl{V}_\Delta = \mcl{V}_\Delta^*$
if the inequality holds for all
$z_\Delta \in L_{\mrm{e}, [a,b]}$ and all $T \geq 0$,
\begin{equation*}
\ip{p_T \Psi_\Delta \bmat{z_\Delta\\ \Delta z_\Delta}}
{p_T \mcl{M}_{\mcl{V}_\Delta}\Psi_\Delta\bmat{z_\Delta\\ \Delta z_\Delta}}
_{L_{[a,b]}}
\geq 0.
\end{equation*}
Moreover, let
$\mcl{V}_\Delta\{\mcl{Q}_\Delta, \mcl{S}_\Delta, \mcl{R}_\Delta\}$
be parameterized as,
\begin{equation*}
\mcl{V}_\Delta =
\bmat{
\mcl{Q}_\Delta & \mcl{S}_\Delta \\
\mcl{S}_\Delta^* & \mcl{R}_\Delta
}
\end{equation*}
\end{definition}
Beyond stability, we are interested in ensuring that the system meets certain
performance specifications in the presence of uncertainty.
Robust performance extends the notion of robust stability by additionally
requiring that a performance metric is satisfied. The associated
input--output maps are evaluated from the zero initial condition
in~(\ref{ass:initial-conditions}).
\begin{definition}
\label{def:robust-performance}\emph{Robust performance} is achieved if the interconnection $P \star \Delta$
is {robustly stable} and for performance IQC $(\Psi_{\mrm{p}}, {\mcl V_{\mrm{p}}})$ defined by $\Psi_{\mrm{p}}:
L^{\mbf{n}_{\mrm{z}_{\mrm{p}}} + \mbf{n}_{\mrm{w}_{\mrm{p}}}}_{\mrm{e}, [a,b]}
\to
L^{\mbf{n}_{\tilde{\mrm{z}}_{\mrm{p}}} + \mbf{n}_{\tilde{\mrm{w}}_{\mrm{p}}}}_{\mrm{e}, [a,b]}$ and
$\mcl{V}_{\mrm{p}}\{\mcl{Q}_\mrm{p}, \mcl{S}_\mrm{p}, \mcl{R}_\mrm{p}\}=\mcl{V}_{\mrm{p}}^*$, such that
\[\ip{p_T \Psi_{\mrm{p}}\bmat{v \\ 0}}{p_T\mcl{M}_{\mcl V_{\mrm{p}}}\Psi_{\mrm{p}}\bmat{v\\ 0}}_{L_{[a,b]}}\geq 0\]
for any $v\in L_{\mrm{e},[a,b]}^{\mbf{n}_{\mrm{z}_{\mrm{p}}}}$ we have 
\begin{gather*}
\ip{p_T\Psi_{\mrm{p}}\bmat{P\star\Delta\\I}w_{\mrm{p}}}{p_T\mcl{M}_{\mcl{V}_{\mrm{p}}}\Psi_{\mrm{p}}\bmat{P\star\Delta\\I}w_{\mrm{p}}}_{L_{[a,b]}}
\\ \leq -\epsilon\|p_Tw_{\mrm{p}}\|_{L_{[a,b]}}^2.
\end{gather*}
% with 
% \begin{equation*}
% \ip{p_T \Psi_{\mathrm{p}}\bmat{v\\ 0}}{p_T\mcl{M}_{\mcl V_{\mathrm{p}}}\Psi_{\mathrm{p}}\bmat{v \\ 0}}_{\mbf L_{[a,b]}}
% \geq 0.
% \end{equation*}
for all $w_{\mrm{p}}\in L_{\mrm{e},[a,b]}^{\mbf{n}_{\mrm{w}_{\mrm{p}}}}$ and $T \geq 0$. 
\end{definition}
In this framework, the multipliers $\mcl M_{\mcl V _\Delta}$ and
$\mcl M_{{\mcl{V}}_{\mrm{p}}}$ serve distinct but complementary roles.
The hard uncertainty IQC $(\Psi_\Delta, \mcl M_{\mcl V_\Delta})$ characterizes the behavior of uncertainties,
nonlinearities, and other perturbations captured by the operator $\Delta$,
ensuring that the IQC is satisfied for the class of admissible disturbances.
In contrast, the performance IQC $(\Psi_{\mrm{p}}, \mcl M_{\mcl V_{\mrm{p}}})$ encodes performance
requirements for the nominal system, such as $L_2$-gain, passivity
conditions, or frequency weighted input-output properties.
Together, these multipliers enable simultaneous verification of robust stability
against perturbations and achievement of specified performance objectives.
% The interconnected structure for robust stability and performance analysis
% can be viewed in Figure~\ref{fig:block structure}.
This section applies the IQC framework to the primal system and gives a
certificate for simultaneous robust stability and performance. For brevity,
we write the hard IQC as
\begin{equation*}
    \ip{p_T \Psi\bmat{Pw \\ w}}{p_T \mcl{M}_\mcl{V}\Psi\bmat{Pw \\ w}}_{L_{[a,b]}} \leq -\epsilon \|p_T w\|^2_{L_{[a,b]}}.
\end{equation*}
Let $\Psi_\Delta$ and $\Psi_{\mrm{p}}$ have a $2\times2$ block partition and define
$\Psi_{ij}:=\operatorname{diag}((\Psi_\Delta)_{ij},(\Psi_{\mrm{p}})_{ij})$ for
$i,j\in\{1,2\}$. Let $w^\top=\bmat{w_\Delta^\top & w_{\mrm{p}}^\top}$ and $z^\top=\bmat{z_\Delta^\top & z_{\mrm{p}}^\top}$.
Moreover, define $\mcl{Q} = \operatorname{diag}(\mcl{Q}_\Delta, \mcl{Q}_{\mrm{p}})$,
$\mcl{S} = \operatorname{diag}(\mcl{S}_\Delta, \mcl{S}_{\mrm{p}})$, and
$\mcl{R} = \operatorname{diag}(\mcl{R}_\Delta, \mcl{R}_{\mrm{p}})$, then
$\mcl{V}\{\mcl{Q}, \mcl{S}, \mcl{R}\}$.

\begin{theorem}[Robust Stability and Performance]\label{thm:robust-stability-performance}
Let $P$, $\mbf{\Delta}$ be given, and there exists $\Psi_\Delta$,
${\mcl{V}_\Delta}$, such that,
\begin{itemize}
\item The interconnection $P\star\Delta$ is well--posed for all
$\Delta\in\mbf{\Delta}$.
\item Any uncertainty
$\Delta\in\mbf{\Delta}$ satisfy the hard IQC defined by $\Psi_\Delta$ and $\mcl V_\Delta$ .
\end{itemize}
Then the interconnection $P \star \Delta$ is \emph{robustly stable} and
\emph{robust performance} is achieved with respect to some $\Psi_{\mrm{p}}$ and
${{\mcl{V}}_{\mrm{p}}} =\{\mcl Q_{\mrm{p}}, \mcl S_{\mrm{p}}, \mcl R_{\mrm{p}}\}$ if,
\begin{equation}\label{eq:robust-condition}
\ip{p_T\Psi\bmat{Pw \\ w}}{p_T\mcl{M}_{\mcl V}\Psi\bmat{Pw \\ w}}_{L_{[a,b]}} 
\leq
-\epsilon
\|p_Tw\|^2_{L_{[a,b]}}
\end{equation}
for all $w_\Delta \in L_{\mrm{e},[a,b]}^{\mbf{n}_{\mrm{w}_\Delta}}$,
$w_{\mrm{p}} \in L_{\mrm{e},[a,b]}^{\mbf{n}_{\mrm{w}_{\mrm{p}}}}$ and $T \geq 0$.
\end{theorem}
\begin{proof}
See \ref{proof:robust-stability-performance}.
\end{proof}
Theorem~\ref{thm:robust-stability-performance} establishes a framework for deriving Linear PI
Inequalities (LPIs), which are used in assessing the robust performance of a
system. We next specialize this certificate to the controlled primal
realization and expose bilinearities in the LPI that distinguish convex analysis from
convex controller synthesis.
\primalsystemcomponents